\documentclass[11pt,letterpaper]{article}
\usepackage[utf8]{inputenc}
\usepackage[T1]{fontenc}
\usepackage[spanish,es-noshorthands,es-tabla]{babel}
\usepackage[margin=2cm,top=2.5cm,bottom=2.5cm]{geometry}
\usepackage{amsmath,amssymb}
\usepackage{enumitem}
\usepackage{fancyhdr}
\usepackage{listings}
\usepackage{xcolor}
\usepackage{tcolorbox}
\usepackage{booktabs}
\usepackage{array}
\raggedbottom
\setlength{\headheight}{14pt}
\let\vspaceoriginal\vspace
\renewcommand{\vspace}[1]{\par\vspaceoriginal{#1}\par}
\pagestyle{fancy}
\fancyhf{}
\fancyhead[L]{Basic C Without Tears}
\fancyhead[C]{\NombreDocumento}
\fancyhead[R]{Stack}
\fancyfoot[L]{Basic-C-Without-Tears}
\fancyfoot[R]{\thepage}
\renewcommand{\headrulewidth}{0.4pt}
\renewcommand{\footrulewidth}{0.4pt}
\lstdefinestyle{customc}{
  language=C,
  numbers=left,
  stepnumber=1,
  numbersep=8pt,
  numberstyle=\tiny\color{gray},
  basicstyle=\footnotesize\ttfamily,
  keywordstyle=\bfseries\color{blue!80!black},
  commentstyle=\itshape\color{green!50!black},
  stringstyle=\color{red!70!black},
  breaklines=true,
  breakatwhitespace=true,
  tabsize=4,
  showstringspaces=false,
  frame=none
}
\newtcolorbox{solucionbox}{
  before=\par,
  after=\par,
  colback=white,
  colframe=black,
  arc=0mm,
  boxrule=0.6pt,
  left=3mm,
  right=3mm,
  top=2mm,
  bottom=2mm
}
\newif\ifpauta
\ifdefined\PAUTA\pautatrue\else\pautafalse\fi

\newcommand{\NombreDocumento}{Examen 03 -- Variante A}
\begin{document}
\begin{center}
    {\LARGE \textbf{Basic C Without Tears}} \\[0.2cm]
    {\Large \textbf{Examen 03 -- Variante A}} \\[0.12cm]
    \small Fundamentos de C, matrices, control de flujo y cadenas de caracteres \\[0.08cm]
    \textbf{Autor:} Stack
\end{center}
\vspace{0.2cm}
\begin{enumerate}[label=\textbf{\arabic*.}]
\item Una matriz contiene las ventas de dos productos durante tres días:
\begin{lstlisting}[style=customc, numbers=none]
int V[3][2] = {{5, 8}, {7, 4}, {6, 9}};
\end{lstlisting}
\begin{enumerate}[label=(\alph*)]
\item ¿Cuántas posiciones hay? Indique el valor de \texttt{V[2][0]}.
\ifpauta\begin{solucionbox}Hay $3\times2=6$ posiciones. \texttt{V[2][0]} vale 6.\end{solucionbox}\else\vspace{1.5cm}\fi
\item Calcule el total vendido por producto (cada columna).
\ifpauta\begin{solucionbox}Producto 1: $5+7+6=18$. Producto 2: $8+4+9=21$.\end{solucionbox}\else\vspace{2cm}\fi
\item Escriba código que cuente cuántas ventas de la matriz son mayores o iguales a 7.
\ifpauta\begin{solucionbox}\begin{lstlisting}[style=customc]
int cantidad = 0;
for (int i = 0; i < 3; i++) {
    for (int j = 0; j < 2; j++) {
        if (V[i][j] >= 7) {
            cantidad++;
        }
    }
}
printf("Ventas de 7 o mas: %d\n", cantidad);
\end{lstlisting}Los valores que cumplen son 8, 7 y 9; el resultado es 3.\end{solucionbox}\else\vspace{3cm}\fi
\end{enumerate}
\newpage
\item Una cantidad ingresada en la variable \texttt{int cantidad} debe ser positiva y no superar 50.
\begin{enumerate}[label=(\alph*)]
\item Escriba una condición que identifique los datos inválidos.
\ifpauta\begin{solucionbox}\texttt{cantidad <= 0 || cantidad > 50}.\end{solucionbox}\else\vspace{1.5cm}\fi
\item Escriba un fragmento que imprima \texttt{Dato aceptado} solo si el valor es válido; en caso contrario, imprima \texttt{Revisar dato}.
\ifpauta\begin{solucionbox}\begin{lstlisting}[style=customc]
if (cantidad > 0 && cantidad <= 50) {
    printf("Dato aceptado\n");
} else {
    printf("Revisar dato\n");
}
\end{lstlisting}\end{solucionbox}\else\vspace{2cm}\fi
\item Clasifique los valores \texttt{0}, \texttt{1}, \texttt{50} y \texttt{51}.
\ifpauta\begin{solucionbox}0: inválido; 1: válido; 50: válido; 51: inválido.\end{solucionbox}\else\vspace{2cm}\fi
\end{enumerate}
\newpage
\item Determine la salida exacta del código.
\begin{lstlisting}[style=customc]
int A[2][2] = {{3, 1}, {2, 4}};
int mayor = A[0][0];
for (int i = 0; i < 2; i++) {
    for (int j = 0; j < 2; j++) {
        if (A[i][j] > mayor) {
            mayor = A[i][j];
        }
    }
}
printf("Mayor=%d\n", mayor);
\end{lstlisting}
\ifpauta\begin{solucionbox}Se comparan sucesivamente 3, 1, 2 y 4 con el mayor actual. El máximo final es 4.\par\textbf{Salida:}\texttt{Mayor=4}.\end{solucionbox}\else\vspace{3cm}\fi
\item Escriba una función que copie una cadena a otra, reemplazando cada espacio por un guion bajo. Suponga que el arreglo de destino tiene espacio suficiente y termine correctamente la cadena.
\begin{center}\begin{tabular}{ll}\toprule\textbf{Entrada} & \textbf{Resultado}\\\midrule\texttt{"hola c"} & \texttt{"hola\_c"}\\\texttt{"sin espacios"} & \texttt{"sin\_espacios"}\\\texttt{"C"} & \texttt{"C"}\\\bottomrule\end{tabular}\end{center}
\ifpauta\begin{solucionbox}\begin{lstlisting}[style=customc]
void copiar_con_guion(const char origen[], char destino[]) {
    int i = 0;
    while (origen[i] != '\0') {
        if (origen[i] == ' ') {
            destino[i] = '_';
        } else {
            destino[i] = origen[i];
        }
        i++;
    }
    destino[i] = '\0';
}
\end{lstlisting}El índice de escritura avanza en paralelo al de lectura y se copia el terminador.\end{solucionbox}\else\vspace{3cm}\fi
\end{enumerate}
\end{document}
